\subsection{Forced central involutions and correlated semion traces}
\label{subsec:frontier-quartic-semion}

The complete integer Gram and distinguished Verlinde enumerations for
$R_{8,3,928}$ and $R_{8,3,941}$ each leave $343$ restriction patterns.
Their previously verified Galois--pseudo-unitary gates permit the following
necessary conditions to be used for general complex categorifiability.

\begin{proposition}
\label{prop:frontier-forced-quartic-involution}
Every categorification of either quartic ring, after this Galois transport,
has a unique central invertible $G$ whose restriction is the nontrivial
central source invertible $g$. It satisfies $G^2=\mathbf1$. For
$R_{8,3,928}$ one has $\theta_G^2=-1$, whereas for $R_{8,3,941}$ one has
$\theta_G^2=1$.
\end{proposition}
\begin{proof}
The complete distinguished restriction list contains only the tensor unit
with restriction supported on $\mathbf1$. Adjunction therefore makes the
forgetful kernel $\mathrm{Vec}$. If a central simple restricts to $m g$,
its product with its dual restricts to $m^2\mathbf1$; the kernel statement
and rigidity force $m=1$. Two such lifts differ by a kernel invertible,
so there is at most one. All $343$ complete integer patterns force exactly
one such lift. The same kernel argument gives $G^2=\mathbf1$.

The complete earlier rotation certificates force
$\nu_2(g)=-1$ and $+1$, respectively. The forgetful functor preserves the
canonical positive pivotal structures. On the pointed order-two
subcategory, the indicator is the normalized cocycle value
$\omega(g,g,g)\in\{1,-1\}$; this also follows from
\cite[Definition~2.2 and Proposition~4.1]{SchauenburgPointedFrontier}
with trivial subgroup. The order-two hexagon gives
$c(G,G)^2=\omega(g,g,g)$, and balancing gives
$1=\theta_G^2c(G,G)^2$. Hence $\theta_G^2=\nu_2(g)$.
\end{proof}

For the semionic case, the monodromy charge of any central simple $Z$
is $\pm1$, since it is an eigenvalue of the involution fusion matrix
$N_G$. Balancing gives
\[
 \theta_{GZ}=\theta_G M_{G,Z}\theta_Z.
\]
Thus $G$ has no fixed simple. The $40$ high restriction vectors fixed by
multiplication by $g$ must have even multiplicity; exactly $15$ of the
$343$ integer patterns remain. The five distinguished low objects
translate to five distinct high objects of multiplicity one, with twists
$\pm i$. Their normalized relative moments at powers $1,\ldots,6$ are
$(0,-1,0,1,0,-1)$, and the second and sixth total moments of every fixed
restriction type vanish.

Imposing these conditions together with the full height-six tensor
rotation equations, all distinguished $STS$ equations, and the fixed
integer Gram pattern gives exact Farkas contradictions for five patterns
and exact rational feasible vectors for ten. This calculation uses the
matrix-character normalization $D/(f_E\dim E)$ throughout.

The remaining correlation can be retained directly. Write
$K=\mathbb Q(\sqrt2,\sqrt3)$ and let $\operatorname{tr}_K$ denote normalized
relative trace. Each unordered $G$-orbit contributes a joint tuple
\[
 \bigl(\operatorname{tr}_K(\theta^k+(i\theta)^k)\bigr)_{k=1}^{6}.
\]
The opposite sign is covered by interchanging the two roots.
Proposition~\ref{prop:m3-biquadratic-traces} gives the finite order box
$288$, together with the all-zero tail. Multiplication by $i$ cannot
remove a prime-power valuation exceeding this bound. Direct exact
averaging of all powers of $\zeta_{288}$ over the $24$ automorphisms fixing
$K$ produces exactly $30$ joint pair types. A pair variable has
multiplicity mass two. Replacing each fixed restriction block by these
pair columns preserves every Gram, rotation, moment, and $STS$ row;
the other high columns and auxiliary rotation variables are retained once.
All ten patterns still have exact rational feasible vectors in the
resulting $873$-equation, $2227$-variable systems.

The simultaneous eigenvalue tuple
\[
 \left(\frac{D\,\chi_E(FZ)}{f_E\dim(E)\dim(Z)}\right)_E
\]
has Galois-invariant multiplicity: the distinguished center fusion
matrices are commuting integer matrices, and Galois transport preserves
the dimensions of their joint eigenspaces. Testing all four automorphisms
of $K$ leaves $255$ of the original $343$ patterns for each quartic ring;
all $15$ semion-parity patterns pass this additional condition.
The same full-tuple test preserves all four complete integer patterns for
$R_{8,3,721}$ and all $1238$ cubic/spectral integer patterns for
$R_{8,3,1154}$, closing this finite test across the four remaining
noncommutative cases. An actual
$\mathcal Z(\mathrm{Fib})$ control checks the joint-spectrum normalization.

The pair variables themselves count actual orbits and hence are integers;
normalized relative trace changes their values, not their multiplicities.
The auxiliary rotation variables similarly count individual eigenvalues.
A bounded integer search finds exact nonnegative integer witnesses for
four of the ten remaining patterns, verified against all $873$ equations.
The other six integer-feasibility questions remain uncertified: numerical
infeasibility is not used as an obstruction. Exact modular lattice tests
at $2,3,5,7,11,13$, including removal of every variable forced zero by
nonnegative rows, do not resolve them.

These vectors are feasible necessary data, not constructions. In
particular, the forced central semion does not imply a tensor
factorization of the original category. The source-bound read-only
entry point for both finite campaigns is
\begin{quote}\small
\texttt{r928\_semion\_rotation/verify\_all.py}.
\end{quote}
The independent orbit and correlated-trace checkers are registered
separately. Actual double-semion center data control the $STS$ and twist
normalizations.
